\section{Unités: difficultés et contradictions}
Au cours d'un entretien très fructueux, Mme Bonzi me dit qu'elle a du mal à faire mettre les unités à ses élèves de terminale scientifique, option physique. Je lui explique ce que je pense en être la cause: au primaire elles sont interdites sous prétexte qu'en mathématiques il n'y aurait pas d'unités. Pourtant, on fait de la physique si on parle de km/h! Mme Bonzi semble être bien d'accord. 

J'en discute à la maison et j'apprends d'Augustin qu'au secondaire dans les petites classes c'est interdit, jusqu'en 4\ieme{} probablement, et \textbf{cela même en cours de physique}! Il se souvient avoir demandé à M. Debussy pourquoi (car je lui ai toujours dit de les mettre) et M. Debussy a finalement accepté de voir les unités. Pendant cette discussion, Jeanne renchérit: les unités sont interdites, alors on enlève des points à ceux qui les mettent, puis elles deviennent obligatoires, alors on enlève des points quand on les oublie.
